\chapter{Rock-Paper-Scissors}

\section{A Small Game, A Serious Battle}
Nikos and Eleni are two children in the school yard. They play
rock-paper-scissors every day.

At first, they laugh and play for fun. But soon the game becomes a real
competition. Nikos wants to prove he can read Eleni's mind. Eleni wants to show
that she can never be predicted. Both of them want to win, not just once, but
again and again.

That is what makes this tiny game interesting: the rules are simple, but the
battle is psychological. If your opponent can guess your next move, you lose.
If you can stay unpredictable, you survive.

\section{Rules of the Game}
Each player chooses one of three symbols at the same time:
\begin{itemize}
  \item Rock
  \item Paper
  \item Scissors
\end{itemize}

The winner is decided by these three relations:
\begin{itemize}
  \item Rock beats scissors.
  \item Scissors beats paper.
  \item Paper beats rock.
\end{itemize}

If both players choose the same symbol, the round is a draw.

So there are only three possible outcomes in each round:
\begin{itemize}
  \item Player A wins
  \item Player B wins
  \item Draw
\end{itemize}

\section{Nikos always plays rock}
Nikos decides to use a fixed strategy: he will always choose rock.

At first, this feels powerful to him. When Eleni picks scissors, he wins. When
she picks rock, the round is a draw. Only paper can beat him.

In the first few rounds, Nikos gets some quick wins and becomes confident. But
Eleni watches carefully. After enough rounds, she recognizes the pattern:
Nikos never changes.

Once the pattern is visible, Nikos is in trouble. Eleni can simply choose paper
more often and win again and again. The strategy that looked strong becomes easy
to exploit because it is predictable.

\section{Second stage: a 3-round cycle}
Nikos now changes strategy. Every three rounds, he plays:
\[
\text{rock} \rightarrow \text{scissors} \rightarrow \text{paper} \rightarrow \text{rock} \rightarrow \cdots
\]

At first glance, this looks much better than repeating only rock. The moves are
mixed, and Eleni cannot exploit him instantly.

However, this strategy is still fully predictable once the cycle is discovered.
After observing a few rounds, Eleni can infer Nikos's next move:
\begin{itemize}
  \item if this round is rock, next is scissors;
  \item if this round is scissors, next is paper;
  \item if this round is paper, next is rock.
\end{itemize}

So Eleni can prepare the winning counter one step ahead. This second stage is
more sophisticated than stage one, but it is still a fixed pattern.

\section{Third stage: shuffled triples}
Nikos improves again. He now uses this rule:
\begin{quote}
In every block of three rounds, I will play exactly one rock, one scissors, and
one paper, but the order can change from block to block.
\end{quote}

This sounds much harder to read. There is no single repeated sequence like
``rock, scissors, paper'' anymore.

Still, Eleni can exploit the structure \emph{inside each 3-round block}:
\begin{itemize}
  \item She watches rounds 1 and 2 of the block.
  \item Since each symbol must appear exactly once in the block, she knows the
  only symbol Nikos has not used yet.
  \item In round 3, she plays the symbol that beats that missing one.
\end{itemize}

Example:
\begin{itemize}
  \item Round 1: Nikos plays rock.
  \item Round 2: Nikos plays paper.
  \item Then round 3 must be scissors.
  \item So Eleni plays rock in round 3 and wins.
\end{itemize}

So Nikos has removed one pattern (fixed order), but kept another constraint
(exactly one of each symbol per three rounds). Eleni uses that constraint as
information and turns it into an advantage.

\section{Conclusion: Unpredictability is the key}
The story teaches one central lesson: in repeated rock-paper-scissors, being
truly unpredictable is the key.

A practical way for Nikos is to randomize in advance. For example, he can throw
dice privately and write his future moves on a paper:
\begin{itemize}
  \item \(1,2 \rightarrow\) rock
  \item \(3,4 \rightarrow\) scissors
  \item \(5,6 \rightarrow\) paper
\end{itemize}
As long as Eleni has never seen that paper, she cannot exploit a pattern in
real time.

When both players use this kind of independent randomization, they reach a
stable strategic point called a \textbf{Nash equilibrium}: no player can improve
their expected result by changing strategy alone.

In this equilibrium, each move is played with probability \(1/3\), so for each
player the long-run outcomes are:
\begin{itemize}
  \item win with probability \(1/3\),
  \item lose with probability \(1/3\),
  \item draw with probability \(1/3\).
\end{itemize}
At that point, the game is fair and neither side can gain a systematic edge.

\section{A Real Story: Choosing an Auction House}
Rock-paper-scissors is not only a children's game. It has been used in a famous
high-stakes business decision.

In 2005, Japanese executive Takashi Hashiyama wanted to sell a valuable art
collection (including works by C\'ezanne, Picasso, and Van Gogh). Two major
auction houses, Christie's and Sotheby's, both wanted the sale. Since he viewed
both proposals as strong, he asked them to settle the choice by playing
rock-paper-scissors.

Christie's reportedly asked advice from an 11-year-old girl, who suggested
scissors. Sotheby's chose paper. Scissors beat paper, so Christie's won the
right to run the auction.

This story is a perfect example of the chapter's theme: even in serious
financial settings, people sometimes rely on simple game-theoretic mechanisms
to make a fair, decisive choice \cite{wiki:rock-paper-scissors}.
